\begin{center}
  \large\textbf{ABSTRAK}\\[2ex]
  \large\textbf{\tatitle{}}
\end{center}

\addcontentsline{toc}{chapter}{ABSTRAK}

\vspace{2ex}

\begingroup
% Menghilangkan padding
\setlength{\tabcolsep}{0pt}

\noindent
\begin{tabularx}{\textwidth}{l >{\centering}m{2em} X}
  Nama Mahasiswa / NRP   & : & \name{} / \nrp        \\

  Department & : & \department{} - Fakultas \faculty      \\

  Dosen Pembimbing        & : & 1. \advisor{}   \\
                    &   & 2. \coadvisor{} \\
\end{tabularx}
\endgroup

% Ubah paragraf berikut dengan abstrak dari tugas akhir
Konjungtivitis adalah penyakit mata menular yang membutuhkan diagnosis cepat, namun diagnosis manual seringkali subjektif. Penelitian ini mengusulkan sistem deteksi otomatis menggunakan \textit{Convolutional Neural Network} (CNN) berbasis mekanisme \textit{Attention}. Untuk meningkatkan akurasi dan mengurangi \textit{noise} (seperti kelopak mata atau bulu mata), sistem ini akan melakukan segmentasi pada area sklera terlebih dahulu dan memperkuat fitur vaskular (kemerahan pembuluh darah). Menggunakan dataset publik berjumlah 1.357 citra yang bersifat \textit{imbalanced}, model akan dilatih menggunakan teknik augmentasi atau \textit{cost-sensitive learning}. Diharapkan penggunaan \textit{Attention-driven CNN} dipadukan dengan segmentasi sklera dapat menghasilkan metrik evaluasi (\textit{F1-Score} dan Akurasi) yang tinggi dan \textit{robust}.

% Ubah kata-kata berikut dengan kata kunci dari tugas akhir
\textbf{Kata Kunci: \keywords{}}
